\documentclass[../../main.tex]{subfiles}
\begin{document}
\section{流}\label{section:ch09-流}

下面是另一种可视化与向量场相关的积分曲线族的方式。设 $M$ 是光滑流形，$V\in\mathfrak{X}(M)$，假设对每个点 $p\in M$，$V$ 有唯一的、从 $p$ 出发且对所有 $t\in\mathbb{R}$ 都有定义的积分曲线，记为 $\theta^{(p)}\colon\mathbb{R}\to M$。（并非每条积分曲线都对所有 $t$ 有定义，但为说明起见暂时假设如此。）对每个 $t\in\mathbb{R}$，可以定义映射 $\theta_t\colon M\to M$，将每个 $p\in M$ 映为从 $p$ 出发沿积分曲线走时间 $t$ 后到达的点：
\begin{align*}
\theta_t(p)=\theta^{(p)}(t).
\end{align*}
每个映射 $\theta_t$ 将流形沿积分曲线"滑动"时间 $t$。平移引理表明 $t\mapsto\theta^{(p)}(t+s)$ 是从 $q=\theta^{(p)}(s)$ 出发的 $V$ 的积分曲线；由于假设积分曲线的唯一性，$\theta^{(q)}(t)=\theta^{(p)}(t+s)$。将其翻译为关于映射 $\theta_t$ 的陈述，即得
\begin{align*}
\theta_t\circ\theta_s(p)=\theta_{t+s}(p).
\end{align*}
连同由定义成立的 $\theta_0(p)=\theta^{(p)}(0)=p$，这意味着映射 $\theta\colon\mathbb{R}\times M\to M$ 是加法群 $\mathbb{R}$ 在 $M$ 上的作用。

受这些观察启发，我们定义 $M$ 上的\textbf{全局流}（也称为\textbf{单参数群作用}）为 $M$ 上的连续左 $\mathbb{R}$-作用，即连续映射 $\theta\colon\mathbb{R}\times M\to M$，满足对所有 $s,t\in\mathbb{R}$ 和 $p\in M$：
\begin{align}
\theta\bigl(t,\theta(s,p)\bigr)=\theta(t+s,p),\qquad\theta(0,p)=p. \label{align:ch09-02}
\end{align}

给定 $M$ 上的全局流 $\theta$，我们如下定义两类映射：
\begin{itemize}
\item 对每个 $t\in\mathbb{R}$，定义连续映射 $\theta_t\colon M\to M$ 为 $\theta_t(p)=\theta(t;p)$。定义性质\eqref{align:ch09-02}等价于群律
\begin{align}
\theta_t\circ\theta_s=\theta_{t+s},\qquad \theta_0=\operatorname{Id}_M. \label{align:ch09-03}
\end{align}
与任意连续群作用的情形一样，每个映射 $\theta_t\colon M\to M$ 是同胚；若流是光滑的，则 $\theta_t$ 是微分同胚。

\item 对每个 $p\in M$，定义曲线 $\theta^{(p)}\colon\mathbb{R}\to M$ 为 $\theta^{(p)}(t)=\theta(t;p)$。该曲线的像是 $p$ 在群作用下的轨道。
\end{itemize}

下一个命题表明，每个光滑全局流恰好以我们上面描述的方式由某个光滑向量场的积分曲线导出。若 $\theta\colon\mathbb{R}\times M\to M$ 是光滑全局流，对每个 $p\in M$，我们定义切向量 $V_p\in T_pM$ 为
\begin{align*}
V_p=\theta^{(p)\prime}(0).
\end{align*}
对应 $p\mapsto V_p$ 是 $M$ 上的一个（粗）向量场，称为 $\theta$ 的\textbf{无穷小生成元}，理由稍后解释。

\begin{figure}[ht]
\centering
\includegraphics[width=0.8\textwidth]{fig-09-4.png}
\caption{整体流的无穷小生成元}\label{figure:ch09-infinitesimal-generator}
\end{figure}

\begin{proposition}\label{proposition:ch09-07}
设 $\theta\colon\mathbb{R}\times M\to M$ 是光滑流形 $M$ 上的光滑全局流。$\theta$ 的无穷小生成元 $V$ 是 $M$ 上的光滑向量场，每条曲线 $\theta^{(p)}$ 是 $V$ 的积分曲线。
\end{proposition}

\begin{proof}
为证明 $V$ 光滑，由\cref{proposition:ch08-14}只需证明对定义在开子集 $U\subseteq M$ 上的每个光滑实值函数 $f$，$Vf$ 光滑。对任意此类 $f$ 和任意 $p\in U$，注意到
\begin{align*}
Vf(p)=V_pf=\theta^{(p)\prime}(0)f=\left.\frac{\mathrm{d}}{\mathrm{d}t}\right|_{t=0}f\bigl(\theta^{(p)}(t)\bigr)=\left.\frac{\partial}{\partial t}\right|_{(0;p)}f\bigl(\theta(t;p)\bigr).
\end{align*}
由于 $f(\theta(t;p))$ 经复合是 $(t,p)$ 的光滑函数，其关于 $t$ 的偏导数也是。因此 $Vf(p)$ 光滑依赖于 $p$，故 $V$ 光滑。

接下来需证明 $\theta^{(p)}$ 是 $V$ 的积分曲线，即对所有 $p\in M$ 和所有 $t\in\mathbb{R}$，$\theta^{(p)\prime}(t)=V_{\theta^{(p)}(t)}$。任取 $t_0\in\mathbb{R}$，令 $q=\theta^{(p)}(t_0)=\theta_{t_0}(p)$，则需证 $\theta^{(p)\prime}(t_0)=V_q$。由群律，对所有 $t$，
\begin{align}
\theta^{(q)}(t)=\theta_t(q)=\theta_t\bigl(\theta_{t_0}(p)\bigr)=\theta_{t+t_0}(p)=\theta^{(p)}(t+t_0). \label{align:ch09-04}
\end{align}
因此，对定义在 $q$ 邻域内的任意光滑实值函数 $f$，
\begin{align}
V_qf=\theta^{(q)\prime}(0)f=\left.\frac{\mathrm{d}}{\mathrm{d}t}\right|_{t=0}f\bigl(\theta^{(q)}(t)\bigr)
=\left.\frac{\mathrm{d}}{\mathrm{d}t}\right|_{t=0}f\bigl(\theta^{(p)}(t+t_0)\bigr)=\theta^{(p)\prime}(t_0)f, \label{align:ch09-05}
\end{align}
即所证。
\end{proof}

\begin{example}\label{example:ch09-08}[全局流]
\cref{example:ch09-01}中描述的平面上的两个向量场的积分曲线都对所有 $t\in\mathbb{R}$ 有定义，因此它们生成全局流。利用该例的结果，可以显式写出流。
\begin{enumerate}[(a)]
\item $\mathbb{R}^2$ 上 $V=\partial/\partial x$ 的流是映射 $\sigma\colon\mathbb{R}\times\mathbb{R}^2\to\mathbb{R}^2$，由 $\sigma_t(x,y)=(x+t,y)$ 给出。对每个非零 $t\in\mathbb{R}$，$\sigma_t$ 将平面向右（$t>0$）或向左（$t<0$）平移距离 $|t|$。

\item $W=x\partial/\partial y-y\partial/\partial x$ 的流是映射 $\theta\colon\mathbb{R}\times\mathbb{R}^2\to\mathbb{R}^2$，由 $\theta_t(x,y)=(x\cos t-y\sin t,x\sin t+y\cos t)$ 给出。对每个 $t\in\mathbb{R}$，$\theta_t$ 将平面绕原点旋转角度 $t$。
\end{enumerate}
\end{example}

\subsection*{流的基本定理}

我们已经看到每个光滑全局流产生一个光滑向量场，其积分曲线恰为流所定义的曲线。反过来，我们希望能说每个光滑向量场都是某个光滑全局流的无穷小生成元。然而，容易看出情况并非如此，因为存在光滑向量场，其积分曲线并非对所有 $t\in\mathbb{R}$ 都有定义。下面是两个例子。

\begin{example}\label{example:ch09-09}
设 $M=\mathbb{R}^2\smallsetminus\{0\}$，带标准坐标 $(x,y)$，$V$ 是 $M$ 上的向量场 $\partial/\partial x$。从 $(-1,0)\in M$ 出发的 $V$ 的唯一积分曲线是 $\gamma(t)=(t-1,0)$。但在此情形下，$\gamma$ 不能连续延拓过 $t=1$。由于 $M$ 在原点处有"洞"，这在直观上是显然的；严格证明如下：假设 $\widetilde{\gamma}$ 是 $\gamma$ 过 $t=1$ 的任意连续延拓，则当 $t\nearrow 1$ 时 $\gamma(t)\to\widetilde{\gamma}(1)\in\mathbb{R}^2\smallsetminus\{0\}$。但也可将 $\gamma$ 视为映入 $\mathbb{R}^2$ 的映射（与包含映射 $M\hookrightarrow\mathbb{R}^2$ 复合），由公式显然当 $t\nearrow 1$ 时 $\gamma(t)\to(0,0)$。由于 $\mathbb{R}^2$ 中的极限唯一，矛盾。
\end{example}

\begin{example}\label{example:ch09-10}
更微妙的例子，设 $M=\mathbb{R}^2$，$W=x^2\partial/\partial x$。容易验证从 $(1,0)$ 出发的 $W$ 的唯一积分曲线是
\begin{align*}
\gamma(t)=\left(\frac{1}{1-t},0\right).
\end{align*}
该曲线也不能延拓过 $t=1$，因为当 $t\nearrow 1$ 时其 $x$ 坐标无界。
\end{example}

因此，我们给出如下定义。若 $M$ 是流形，则 $M$ 的一个\textbf{流域}是 $\mathbb{R}\times M$ 的开子集 $\mathscr{D}$，满足对每个 $p\in M$，集合 $\mathscr{D}(p)=\{t\in\mathbb{R}:(t;p)\in\mathscr{D}\}$ 是包含 $0$ 的开区间。$M$ 上的一个\textbf{流}是连续映射 $\theta\colon\mathscr{D}\to M$，其中 $\mathscr{D}\subseteq\mathbb{R}\times M$ 是流域，满足如下群律：对所有 $p\in M$，
\begin{align}
\theta(0;p)=p, \label{align:ch09-06}
\end{align}
以及对所有 $s\in\mathscr{D}(p)$ 和 $t\in\mathscr{D}(\theta(s;p))$ 且 $s+t\in\mathscr{D}(p)$，
\begin{align}
\theta\bigl(t,\theta(s;p)\bigr)=\theta(t+s;p). \label{align:ch09-07}
\end{align}

有时称 $\theta$ 为\textbf{局部流}，以区别于前面定义的全局流。也使用"局部单参数群作用"这一不太常用的术语。

\begin{figure}[ht]
\centering
\includegraphics[width=0.8\textwidth]{fig-09-5.png}
\caption{流的定义域}\label{figure:ch09-flow-domain}
\end{figure}

若 $\theta$ 是流，当 $(t;p)\in\mathscr{D}$ 时，与全局流一样定义 $\theta_t(p)=\theta^{(p)}(t)=\theta(t;p)$。对每个 $t\in\mathbb{R}$，还定义
\begin{align}
M_t=\{p\in M:(t;p)\in\mathscr{D}\}, \label{align:ch09-08}
\end{align}
于是 $p\in M_t\iff t\in\mathscr{D}(p)\iff(t;p)\in\mathscr{D}$。若 $\theta$ 光滑，其无穷小生成元定义为 $V_p=\theta^{(p)\prime}(0)$。

\begin{proposition}\label{proposition:ch09-11}
若 $\theta\colon\mathscr{D}\to M$ 是光滑流，则 $\theta$ 的无穷小生成元 $V$ 是光滑向量场，每条曲线 $\theta^{(p)}$ 是 $V$ 的积分曲线。
\end{proposition}

\begin{proof}
证明与全局流的相应证明\cref{proposition:ch09-07}本质相同。在证明 $V$ 光滑时，只需注意对任意 $p_0\in M$，由于 $\mathscr{D}$ 是开的，$\theta(t;p)$ 对充分接近 $(0;p_0)$ 的所有 $(t;p)$ 都有定义且光滑。在证明 $\theta^{(p)}$ 是积分曲线时，需验证\eqref{align:ch09-04}和\eqref{align:ch09-05}中的所有表达式都有意义。假设 $t_0\in\mathscr{D}(p)$。由于 $\mathscr{D}(p)$ 和 $\mathscr{D}(\theta_{t_0}(p))$ 都是包含 $0$ 的开区间，存在正数 $\varepsilon$ 使得当 $|t|<\varepsilon$ 时 $t+t_0\in\mathscr{D}(p)$ 且 $t\in\mathscr{D}(\theta_{t_0}(p))$，然后由流的定义 $\theta_t(\theta_{t_0}(p))=\theta_{t+t_0}(p)$。证明的其余部分与之前相同。
\end{proof}

下一个定理是本节的主要结果。\textbf{极大积分曲线}是不能延拓为更大开区间上的积分曲线的积分曲线；\textbf{极大流}是不能延拓为更大流域上的流的流。

\begin{theorem}\label{theorem:ch09-12}[流的基本定理]
设 $V$ 是光滑流形 $M$ 上的光滑向量场。存在唯一的光滑极大流 $\theta\colon\mathscr{D}\to M$，其无穷小生成元为 $V$。该流具有如下性质：
\begin{enumerate}[(a)]
\item 对每个 $p\in M$，曲线 $\theta^{(p)}\colon\mathscr{D}(p)\to M$ 是从 $p$ 出发的唯一极大积分曲线。
\item 若 $s\in\mathscr{D}(p)$，则 $\mathscr{D}(\theta(s;p))$ 是区间 $\mathscr{D}(p)-s=\{t-s:t\in\mathscr{D}(p)\}$。
\item 对每个 $t\in\mathbb{R}$，集合 $M_t$ 在 $M$ 中开，且 $\theta_t\colon M_t\to M_{-t}$ 是微分同胚，其逆为 $\theta_{-t}$。
\end{enumerate}
\end{theorem}

\begin{proof}
\cref{proposition:ch09-02}表明从每一点 $p\in M$ 出发的积分曲线存在。假设 $\gamma,\widetilde{\gamma}\colon J\to M$ 是定义在同一开区间 $J$ 上的两条 $V$ 的积分曲线，且在某 $t_0\in J$ 处 $\gamma(t_0)=\widetilde{\gamma}(t_0)$。令 $\Gamma$ 为 $J$ 中满足 $\gamma(t)=\widetilde{\gamma}(t)$ 的 $t$ 的集合。显然 $\Gamma\neq\emptyset$，因为由假设 $t_0\in\Gamma$，且由连续性 $\Gamma$ 在 $J$ 中闭。另一方面，假设 $t_1\in\Gamma$。则在点 $p=\gamma(t_1)$ 周围的光滑坐标邻域内，$\gamma$ 和 $\widetilde{\gamma}$ 都是同一 ODE 的具有相同初始条件 $\gamma(t_1)=\widetilde{\gamma}(t_1)=p$ 的解。由\cref{theorem:appD-01}的唯一性部分，$\gamma$ 和 $\widetilde{\gamma}$ 在包含 $t_1$ 的区间上相等，这意味着 $\Gamma$ 在 $J$ 中开。由于 $J$ 连通，$\Gamma=J$，即 $\gamma=\widetilde{\gamma}$ 在整个 $J$ 上。因此，在一点处一致的任意两条积分曲线在其公共定义域上一致。

对每个 $p\in M$，令 $\mathscr{D}(p)$ 为所有包含 $0$ 且定义了从 $p$ 出发的积分曲线的开区间 $J\subseteq\mathbb{R}$ 的并集。定义 $\theta^{(p)}\colon\mathscr{D}(p)\to M$ 为 $\theta^{(p)}(t)=\gamma(t)$，其中 $\gamma$ 是从 $p$ 出发、定义在包含 $0$ 和 $t$ 的开区间上的任意积分曲线。由于由上述论证所有此类积分曲线在 $t$ 处一致，$\theta^{(p)}$ 定义良好，且显然是从 $p$ 出发的唯一极大积分曲线。

现在令 $\mathscr{D}=\{(t;p)\in\mathbb{R}\times M:t\in\mathscr{D}(p)\}$，定义 $\theta\colon\mathscr{D}\to M$ 为 $\theta(t;p)=\theta^{(p)}(t)$。与通常一样，也写 $\theta_t(p)=\theta(t;p)$。由定义，$\theta$ 满足基本定理陈述中的性质 (a)：对每个 $p\in M$，$\theta^{(p)}$ 是从 $p$ 出发的唯一极大积分曲线。为验证群律，固定任意 $p\in M$ 和 $s\in\mathscr{D}(p)$，记 $q=\theta(s;p)=\theta^{(p)}(s)$。由 $\gamma(t)=\theta^{(p)}(t+s)$ 定义的曲线 $\gamma\colon\mathscr{D}(p)-s\to M$ 从 $q$ 出发，平移引理表明 $\gamma$ 是 $V$ 的积分曲线。由 ODE 解的唯一性，$\gamma$ 与 $\theta^{(q)}$ 在其公共定义域上一致，这等价于第二个群律\eqref{align:ch09-07}，而第一个群律\eqref{align:ch09-06}直接由定义得到。由 $\theta^{(q)}$ 的极大性，$\gamma$ 的定义域不能大于 $\mathscr{D}(q)$，这意味着 $\mathscr{D}(p)-s\subseteq\mathscr{D}(q)$。由于 $0\in\mathscr{D}(p)$，这推出 $-s\in\mathscr{D}(q)$，且群律给出 $\theta^{(q)}(-s)=p$。将同样的论证用 $(-s;q)$ 代替 $(s;p)$，得 $\mathscr{D}(q)+s\subseteq\mathscr{D}(p)$，即 $\mathscr{D}(q)\subseteq\mathscr{D}(p)-s$。这证明了 (b)。

接下来证明 $\mathscr{D}$ 在 $\mathbb{R}\times M$ 中开（故为流域），且 $\theta\colon\mathscr{D}\to M$ 光滑。定义 $\mathscr{D}$ 的子集 $W$ 为所有 $(t;p)\in\mathscr{D}$ 的集合，使得 $\theta$ 在形如 $J\times U\subseteq\mathscr{D}$ 的 $(t;p)$ 的乘积邻域上有定义且光滑，其中 $J\subseteq\mathbb{R}$ 是包含 $0$ 和 $t$ 的开区间，$U\subseteq M$ 是 $p$ 的邻域。则 $W$ 在 $\mathbb{R}\times M$ 中开，$\theta$ 在 $W$ 上的限制光滑，故只需证明 $W=\mathscr{D}$。假设不然，则存在某点 $(\tau;p_0)\in\mathscr{D}\smallsetminus W$。为简单起见假设 $\tau>0$；$\tau<0$ 的论证类似。

\begin{figure}[ht]
\centering
\includegraphics[width=0.8\textwidth]{fig-09-6.png}
\caption{$D$ 是开集的证明}\label{figure:ch09-proof-d-open}
\end{figure}

令 $t_0=\inf\{t\in\mathbb{R}:(t;p_0)\notin W\}$。由 ODE 定理（应用于 $p_0$ 周围的光滑坐标），$\theta$ 在 $(0;p_0)$ 的某乘积邻域上有定义且光滑，故 $t_0>0$。由于 $t_0\leqslant\tau$ 且 $\mathscr{D}(p_0)$ 是包含 $0$ 和 $\tau$ 的开区间，故 $t_0\in\mathscr{D}(p_0)$。令 $q_0=\theta^{(p_0)}(t_0)$。再次由 ODE 定理，存在 $\varepsilon>0$ 和 $q_0$ 的邻域 $U_0$ 使得 $(-\varepsilon,\varepsilon)\times U_0\subseteq W$。我们将用群律证明 $\theta$ 光滑延拓到 $(t_0;p_0)$ 的邻域，这与 $t_0$ 的选择矛盾。

选取 $t_1<t_0$ 使得 $t_1+\varepsilon>t_0$ 且 $\theta^{(p_0)}(t_1)\in U_0$。由于 $t_1<t_0$，有 $(t_1;p_0)\in W$，故存在乘积邻域 $(t_1-\delta,t_1+\delta)\times U_1\subseteq W$。由 $W$ 的定义，这意味着 $\theta$ 在 $[0,t_1+\delta)\times U_1$ 上有定义且光滑。由于 $\theta(t_1;p_0)\in U_0$，可将 $U_1$ 取得足够小使得 $\theta$ 将 $\{t_1\}\times U_1$ 映入 $U_0$。定义 $\widetilde{\theta}\colon[0,t_1+\varepsilon)\times U_1\to M$ 为
\begin{align*}
\widetilde{\theta}(t;p)=
\begin{cases}
\theta_t(p), & p\in U_1,\;0\leqslant t<t_1,\\
\theta_{t-t_1}\circ\theta_{t_1}(p), & p\in U_1,\;t_1-\varepsilon<t<t_1+\varepsilon.
\end{cases}
\end{align*}
$\theta$ 的群律保证这些定义在重叠处一致，而对 $U_1,t_1,\varepsilon$ 的选取保证这定义了一个光滑映射。由平移引理，每个映射 $t\mapsto\widetilde{\theta}(t;p)$ 是 $V$ 的积分曲线，故 $\widetilde{\theta}$ 是 $\theta$ 到 $(t_0;p_0)$ 邻域的光滑延拓，与 $t_0$ 的选择矛盾。这完成了 $W=\mathscr{D}$ 的证明。

最后证明 (c)。$M_t$ 开是 $\mathscr{D}$ 开的直接推论。由 (b) 推出
\begin{align*}
p\in M_t\Rightarrow t\in\mathscr{D}(p)\Rightarrow\mathscr{D}(\theta_t(p))=\mathscr{D}(p)-t\Rightarrow -t\in\mathscr{D}(\theta_t(p))\Rightarrow\theta_t(p)\in M_{-t},
\end{align*}
这表明 $\theta_t$ 将 $M_t$ 映入 $M_{-t}$。此外，群律表明 $\theta_{-t}\circ\theta_t$ 在 $M_t$ 上等于恒等映射。交换 $t$ 和 $-t$ 的角色表明 $\theta_t\circ\theta_{-t}$ 在 $M_{-t}$ 上是恒等映射，证明完毕。
\end{proof}

基本定理所断言存在且唯一的流称为由 $V$\textbf{生成的流}，或简称 $V$ 的\textbf{流}。"无穷小生成元"这一术语来自如下图像：在光滑坐标卡中，积分曲线的良好近似可通过将许多小直线运动复合得到，每步运动的方向和长度由上一步到达点处向量场的值决定。直观上，可以把流想象为无穷多个无穷小线性步的序列。

积分曲线的自然性（\cref{proposition:ch09-06}）翻译为如下流的自然性陈述。

\begin{proposition}\label{proposition:ch09-13}[流的自然性]
假设 $M$ 和 $N$ 是光滑流形，$F\colon M\to N$ 是光滑映射，$X\in\mathfrak{X}(M)$，$Y\in\mathfrak{X}(N)$。设 $\theta$ 是 $X$ 的流，$\psi$ 是 $Y$ 的流。若 $X$ 和 $Y$ 是 $F$-相关的，则对每个 $t\in\mathbb{R}$，$F(M_t)\subseteq N_t$ 且在 $M_t$ 上 $\psi_t\circ F=F\circ\theta_t$。
\end{proposition}

\begin{proof}
由\cref{proposition:ch09-06}，对任意 $p\in M$，曲线 $F\circ\theta^{(p)}$ 是从 $F\circ\theta^{(p)}(0)=F(p)$ 出发的 $Y$ 的积分曲线。因此由积分曲线的唯一性，极大积分曲线 $\psi^{(F(p))}$ 至少在区间 $\mathscr{D}(p)$ 上有定义，且在该区间上 $F\circ\theta^{(p)}=\psi^{(F(p))}$。这意味着
\begin{align*}
p\in M_t\Rightarrow t\in\mathscr{D}(p)\Rightarrow t\in\mathscr{D}(F(p))\Rightarrow F(p)\in N_t,
\end{align*}
即 $F(M_t)\subseteq N_t$，且
\begin{align*}
F\bigl(\theta^{(p)}(t)\bigr)=\psi^{(F(p))}(t)\quad\text{对所有 }t\in\mathscr{D}(p),
\end{align*}
即对所有 $p\in M_t$，$\psi_t\circ F(p)=F\circ\theta_t(p)$。
\end{proof}

下一个推论是直接的。

\begin{corollary}\label{corollary:ch09-14}[流的微分同胚不变性]
设 $F\colon M\to N$ 是微分同胚。若 $X\in\mathfrak{X}(M)$ 且 $\theta$ 是 $X$ 的流，则 $F_*X$ 的流为 $\psi_t=F\circ\theta_t\circ F^{-1}$，其定义域为 $N_t=F(M_t)$（对每个 $t\in\mathbb{R}$）。
\end{corollary}

\subsection*{完备向量场}

如本章前面所观察的，并非每个光滑向量场都生成全局流。生成全局流的向量场足够重要，值得命名。我们说光滑向量场是\textbf{完备的}，若它生成全局流，等价地，若其每条极大积分曲线都对所有 $t\in\mathbb{R}$ 有定义。例如，\cref{example:ch09-08}中计算了流的平面上的两个向量场是完备的，而\cref{example:ch09-09}和\cref{example:ch09-10}中的则不是。

仅凭观察向量场并不总是容易判断它是否完备。若能显式求解 ODE 找到所有积分曲线，且它们都对所有时间存在，则向量场完备。另一方面，若能找到一条不能延拓到整个 $\mathbb{R}$ 的积分曲线（如\cref{example:ch09-09}和\cref{example:ch09-10}中的向量场所做的），则它不完备。然而，显式求解 ODE 通常不可能，因此有一些判断向量场何时完备的一般准则是有用的。

下面将证明所有紧支集光滑向量场（因而紧流形上的所有光滑向量场）都是完备的。证明基于如下引理。

\begin{lemma}\label{lemma:ch09-15}[一致时间引理]
设 $V$ 是光滑流形 $M$ 上的光滑向量场，$\theta$ 是其流。假设存在正数 $\varepsilon$ 使得对每个 $p\in M$，$\theta^{(p)}$ 的定义域都包含 $(-\varepsilon,\varepsilon)$。则 $V$ 完备。
\end{lemma}

\begin{proof}
为反证，假设对某 $p\in M$，$\theta^{(p)}$ 的定义域 $\mathscr{D}(p)$ 上有界。（下有界的证明类似。）令 $b=\sup\mathscr{D}(p)$，取正数 $t_0$ 使得 $b-\varepsilon<t_0<b$，令 $q=\theta^{(p)}(t_0)$。假设意味着 $\theta^{(q)}(t)$ 至少对 $t\in(-\varepsilon,\varepsilon)$ 有定义。定义曲线 $\gamma\colon(-\varepsilon,t_0+\varepsilon)\to M$ 为
\begin{align*}
\gamma(t)=
\begin{cases}
\theta^{(p)}(t), & -\varepsilon<t<b,\\
\theta^{(q)}(t-t_0), & t_0-\varepsilon<t<t_0+\varepsilon.
\end{cases}
\end{align*}
这两个定义在重叠处一致，因为由 $\theta$ 的群律 $\theta^{(q)}(t-t_0)=\theta_{t-t_0}(q)=\theta_{t-t_0}\circ\theta_{t_0}(p)=\theta_t(p)=\theta^{(p)}(t)$。由平移引理，$\gamma$ 是从 $p$ 出发的积分曲线。由于 $t_0+\varepsilon>b$，矛盾。
\end{proof}

\begin{theorem}\label{theorem:ch09-16}
光滑流形上的每个紧支集光滑向量场都是完备的。
\end{theorem}

\begin{proof}
设 $V$ 是光滑流形 $M$ 上的紧支集向量场，$K=\operatorname{supp}V$。对每个 $p\in K$，存在 $p$ 的邻域 $U_p$ 和正数 $\varepsilon_p$ 使得 $V$ 的流至少在 $[-\varepsilon_p,\varepsilon_p]\times U_p$ 上有定义。由紧性，有限个这样的集合 $U_{p_1},\dots,U_{p_k}$ 覆盖 $K$。取 $\varepsilon=\min\{\varepsilon_{p_1},\dots,\varepsilon_{p_k}\}$，则从 $K$ 出发的每条极大积分曲线至少在 $(-\varepsilon,\varepsilon)$ 上有定义。由于 $V$ 在 $K$ 外恒为零，从 $M\smallsetminus K$ 出发的每条积分曲线都是常值的，因而可定义在整个 $\mathbb{R}$ 上。故一致时间引理的假设满足，$V$ 完备。
\end{proof}

\begin{corollary}\label{corollary:ch09-17}
在紧光滑流形上，每个光滑向量场都是完备的。
\end{corollary}

Lie 群上的左不变向量场构成另一类总是完备的向量场。

\begin{theorem}\label{theorem:ch09-18}
Lie 群上的每个左不变向量场都是完备的。
\end{theorem}

\begin{proof}
设 $G$ 是 Lie 群，$X\in\operatorname{Lie}(G)$，$\theta\colon\mathscr{D}\to G$ 表示 $X$ 的流。存在某 $\varepsilon>0$ 使得 $\theta^{(e)}$ 在 $(-\varepsilon,\varepsilon)$ 上有定义。

任取 $g\in G$。由于 $X$ 与自身 $L_g$-相关，由\cref{proposition:ch09-06}，曲线 $L_g\circ\theta^{(e)}$ 是从 $g$ 出发的 $X$ 的积分曲线，因而等于 $\theta^{(g)}$。这表明对每个 $g\in G$，积分曲线 $\theta^{(g)}$ 至少在 $(-\varepsilon,\varepsilon)$ 上有定义，故一致时间引理保证 $X$ 完备。
\end{proof}

下面是积分曲线的另一个有用性质。

\begin{lemma}\label{lemma:ch09-19}[逃逸引理]
假设 $M$ 是光滑流形，$V\in\mathfrak{X}(M)$。若 $\gamma\colon J\to M$ 是 $V$ 的极大积分曲线，其定义域 $J$ 有有限上确界 $b$，则对任意 $t_0\in J$，$\gamma([t_0,b))$ 不包含于 $M$ 的任何紧子集。
\end{lemma}

\begin{proof}
习题 9-6。
\end{proof}

\end{document}